\section{Who pays for load forecasts, and how}

Load forecasting is not one job. Different actors, different money.

\subsection{The balancing mechanism --- where errors become invoices}

Every market participant belongs to a balancing group run by a
\emph{Balance Responsible Party} (BRP). Day-ahead, the BRP declares a
schedule: how much its group will consume or produce, hour by hour.
On delivery day, reality deviates. The TSO fixes the gap in real time and
bills the BRP at the \emph{balancing price} --- systematically worse than the
day-ahead price you could have traded at. Forecast error $\to$ imbalance
$\to$ money burned. This mechanism, not fines or prestige, funds every
forecasting desk.

\subsection{The actors}

\begin{itemize}
  \item \textbf{TSO (PSE).} Forecasts national load for system security:
        reserves, grid limits, unit commitment. Not a trading goal.
        Their forecast is public --- that is why it is our benchmark.
  \item \textbf{Retailers (B2C/B2B sellers --- PGE, Tauron, Orlen Energa).}
        Forecast \emph{their own customers'} consumption, buy that amount
        day-ahead. Goal: minimize their portfolio's imbalance cost.
        They do not care about beating PSE's national number; they fight
        their own error. A 1\% MAPE improvement on a big Polish portfolio
        is millions of PLN per year in avoided balancing charges.
  \item \textbf{Industrial consumers.} Same story, own plant load, often
        outsourced to a vendor.
  \item \textbf{Traders and prop desks (Axpo, Statkraft, Shell).} Load is an
        \emph{input}, not the product. National load + wind + solar +
        outages $\to$ price forecast $\to$ positions (day-ahead vs intraday
        vs balancing spreads). Edge comes from having better fundamentals
        than the consensus. If you predict demand better than the market
        (whose anchor is often the TSO's public forecast), you predict
        prices better. Closer to your ``arbitrage'' intuition --- but the
        profit is realized through price positions, never by literally
        betting against the TSO's MAPE.
  \item \textbf{Forecast vendors (Volue, Meteomatics, Yes Energy).} Sell all
        of the above to everyone. Accuracy tables are their marketing.
\end{itemize}

\subsection{So what is OUR goal?}

This repo simulates the retailer/TSO-style desk first: lowest honest
day-ahead MAPE/pinball vs the PSE benchmark, with explanations. Phase 2
(price forecasting) adds the trader's angle, where our load forecast becomes
one of the fundamental inputs.

\subsection{Interview line}

``A retailer's forecast error is an invoice: imbalance settled at worse
prices. A trader's forecast error is a missed position. I built for the
first and I am extending to the second.''
